\section{The complete generalized Paley table}
To avoid confusing the subgroup index with the valency, write
$\operatorname{GP}(q,h)$ for the graph on $\mathbb F_q$ with connection
subgroup $S=\langle g^h\rangle$, where $g$ generates
$\mathbb F_q^\times$, $h\mid q-1$ and $-1\in S$.
Its degree is $\nu=(q-1)/h$. Translations preserve the graph and both
adjacency and reflexive codes. Thus the exact group hypothesis of
\Uref{ULCD.TransitiveDistance} is available for every row below.

\begin{theorem}\label{thm:paley-table}
All twelve rows of source Table~1 have the following parameters.
Every displayed adjacency code is RLCD and none is CLCD.
A star in the last column means that its complete set of minimum words
is exactly the set of adjacency rows.
\begin{center}\small
\begin{tabular}{rrrlll c}\toprule
$q$&$h$&$p$&$(\nu,\lambda,\mu)$&$C_A$&$C_{I+A}$&\\\midrule
25&3&2&$(8,3,2)$&$[25,8,8]$&$[25,17,4]$&$*$\\
49&4&2&$(12,5,2)$&$[49,12,12]$&$[49,37,4]$&$*$\\
81&4&2&$(20,1,6)$&$[81,20,20]$&$[81,61,6]$&$*$\\
81&5&2&$(16,7,2)$&$[81,16,16]$&$[81,65,4]$&$*$\\
121&3&2&$(40,15,12)$&$[121,40,20]$&$[121,81,8]$&\\
121&4&2&$(30,11,6)$&$[121,90,6]$&$[121,31,11]$&\\
121&4&3&$(30,11,6)$&$[121,30,20]$&$[121,91,6]$&\\
121&6&2&$(20,9,2)$&$[121,20,20]$&$[121,101,4]$&$*$\\
243&11&2&$(22,1,2)$&$[243,110,d]$&$[243,133,\delta]$&\\
729&4&2&$(182,55,42)$&$[729,546,d]$&$[729,183,\delta]$&\\
729&7&2&$(104,31,12)$&$[729,104,d]$&$[729,625,\delta]$&\\
729&14&2&$(52,25,2)$&$[729,52,d]$&$[729,677,\delta]$&\\\bottomrule
\end{tabular}
\end{center}
In the last four rows only $d\leq\nu$ and $\delta\leq\nu+1$ are
asserted. Every bracket describes a code over $\mathbb F_p$.
The printed degree $201$ in the $(81,4)$ source row is corrected to $20$.
\end{theorem}
\begin{proof}
The complete finite certificate constructs each field as
$\mathbb F_\ell[X]/(f)$, recording $f$, the primitive element and the
subgroup $S$. It verifies irreducibility of $f$ and all $q-1$ distinct
nonzero powers of the generator. For every nonzero $t$, it counts
$|S\cap(t+S)|$. Translation reduces every pair of distinct vertices
to $(0,t)$, so those counts prove the displayed strongly regular
parameters, not just a random selection of them.

In each stated characteristic $\nu=\mu=0$ and $\lambda=-1$.
The common SRG criterion gives $A^2=-A$, proving reflexive duality.
The complement is orthogonal, but both degrees vanish, since
$q-1=0$ modulo $p$ in every row. The constant-line defect therefore
proves that none is CLCD. The two nonconstant rational eigenvalues are
the roots of $X^2-(\lambda-\mu)X-(\nu-\mu)$; their multiplicities
satisfy $f+g=q-1$ and $\nu+fr+gs=0$. They are integral in every row.
\Uref{ULCD.Rank} now gives all dimensions, with Gaussian elimination
also providing systematic generators for the first eight rows.

For each of the first eight rows the certificate supplies, for both
codes, a word of the displayed distance $D$ and a complete enumeration
of nonzero messages of weight at most
$\lfloor K(D-1)/q\rfloor$, where $K$ is that code's dimension.
Every enumerated word has weight at least $D$.
The translation action just established allows
\Uref{ULCD.TransitiveDistance} to prove all sixteen exact distances.
For the five starred rows a second complete enumeration, with bound
$\lfloor KD/q\rfloor$, finds only adjacency rows at weight $D$.
Their set is translation invariant, so
\Uref{ULCD.MinimumEnumeration} proves that these are all minimum words.
For each of the last four rows, an adjacency row has weight $\nu$ and
a reflexive row has weight $\nu+1$, proving exactly the asserted upper
bounds and no stronger conclusion.
\end{proof}

The finite premise consists of the whole programs
\path{check_paley_codes.py} and
\path{enumerate_sparse_messages.c}, together with their complete
field, basis, witness and enumeration records. The Python driver checks
systematic bases, membership of each witness, the exact binomial message
count and the lower-bound output. The C enumerator visits every
strictly increasing tuple of row positions and every tuple of nonzero
coefficients once. Its binary and ternary arithmetic and boundary
conditions have a separate accepted review against coordinate oracles.
Those checks are required evidence for the finite distance conclusions;
they are not replaced by a reference to the historical Magma output.
The four large rows have no unreported minimum-distance computation.

The source attributes the original generalized Paley computations to
J.~and P.~Limbupasiriporn (its reference [27]), and cites
Seneviratne--Limbupasiriporn [30] for the neighbour formula.
Indeed multiplication by an element of $S$ makes
$|N(1)\cap S|$ the common adjacent-pair count for every such graph.
The source's broader assertion from [27] that $\lambda-\mu$ is odd for
strongly regular cases with $h>2$ is a retained cited observation; our
certificate verifies it in all twelve displayed rows, and neither the
generic SRG argument nor this table purports to classify all generalized
Paley graphs.
